% Generated draft: requires mathematical and source review.
\documentclass[11pt]{article}
\usepackage[margin=25mm]{geometry}
\usepackage{fontspec}
\setmainfont{Latin Modern Roman}
\usepackage{amsmath,amssymb,mathtools,unicode-math}
\setmathfont{Latin Modern Math}
\usepackage{enumitem,needspace,xurl,hyperref}
\setlength{\parindent}{0pt}
\setlength{\parskip}{5pt}
\setlength{\emergencystretch}{3em}
\newcommand{\N}{\mathbb{N}}
\newcommand{\Z}{\mathbb{Z}}
\newcommand{\Q}{\mathbb{Q}}
\newcommand{\R}{\mathbb{R}}
\newcommand{\C}{\mathbb{C}}
\newcommand{\F}{\mathbb{F}}
\newcommand{\A}{\mathbb{A}}
\newcommand{\PP}{\mathbb{P}}
\newcommand{\E}{\mathbb{E}}
\newcommand{\eps}{\varepsilon}
\newcommand{\dd}{\,\mathrm{d}}
\begin{document}
\section*{UnsolvedMath ID: 30002766}
\textbf{Source problem code:} catalogue entry
\subsection{Definitions and mathematical statement}
For this draft, interpret transitivity degree using faithful permutation actions on a countably infinite set. This is a conventional interpretation; the supplied record does not explicitly define the term, so its agreement with the cited report requires verification.

Let \(G\) be a countably infinite group and let \(X\) be a countably infinite set. A faithful action is an injective homomorphism \(\rho:G\to\operatorname{Sym}(X)\), where \(\operatorname{Sym}(X)\) is the group of all permutations of \(X\). For a positive integer \(m\), put
\[
X^{(m)}=\{(x_1,\ldots,x_m)\in X^m:x_i\ne x_j\text{ whenever }i\ne j\}.
\]
The action is \(m\)-transitive if, for all \((x_1,\ldots,x_m),(y_1,\ldots,y_m)\in X^{(m)}\), there exists \(g\in G\) such that
\[
\rho(g)(x_i)=y_i\qquad(1\le i\le m).
\]
Define
\[
\operatorname{td}(G)=\sup\bigl\{m\in\mathbb Z_{\ge1}:G\text{ admits a faithful }m\text{-transitive action on a countably infinite set}\bigr\},
\]
with value \(\infty\) when the set of such integers is unbounded. Call \(G\) highly transitive if it admits a single faithful action on a countably infinite set that is \(m\)-transitive for every positive integer \(m\). This definition requires one action to work for all \(m\); the definition of \(\operatorname{td}(G)=\infty\) allows the action to depend on \(m\).

\textbf{Existence question.} Under the interpretation above, is the following assertion true?
\[
\forall k\in\mathbb Z_{\ge1}\quad\exists\text{ a countably infinite group }G_k
\quad\operatorname{td}(G_k)=k.
\]
Thus the requested degree is exactly \(k\), rather than merely at least \(k\). The supplied statement writes \(k\in\mathbb N\) without specifying whether zero is included; the displayed formulation explicitly uses the positive-integer interpretation, and the intended convention remains to be checked.

\textbf{Broader classification question.} Is there a reasonable classification of highly transitive groups, or of groups of high transitivity degree? The record does not specify what constitutes a reasonable classification, what threshold is meant by high degree, or whether the countably infinite restriction also governs this broader question. Consequently this part remains a research question without a uniquely determined formal target.
\subsection{Short English statement}
For each positive integer \(k\), ask whether some countably infinite group has transitivity degree exactly \(k\), interpreted through faithful actions on countably infinite sets. The source also asks for a classification of highly transitive groups or groups of high transitivity degree, but leaves the scope and classification criteria unspecified. The supplied open-status assessment is dated curation, not independently verified here.
\subsection{Sources}
\begin{itemize}
\item Local UnsolvedMath registry: supplied complete record for permanent ID \(30002766\), title \emph{Countable Groups of Prescribed Transitivity Degree}, exact source problem code \texttt{catalogue entry}; collection: \emph{Oberwolfach Reports Open Problems}. The record gives proposed year 2015 and supplies no named proposer.
\item Supplied report citation: \emph{Geometric Topology} (2015), Oberwolfach Reports, pp.~187--233, DOI \texttt{10.4171/owr/2015/3}. Literature-source label: ``Osin, Oberwolfach Report.'' \url{https://doi.org/10.4171/owr/2015/3}. The supplied background describes these reports as documenting mathematical workshops and their open-problem discussions.
\item Additional supplied literature source: ``Osin, primary report text.'' \url{https://ems.press/content/serial-article-files/46555}. The supplied attribution to Osin is retained as a reference label; it does not establish a named proposer for this problem.
\item Supplied literature assessment, checked on \textbf{2026-08-21}: status \emph{open}; the record states that existence of countable groups of every finite transitivity degree remains unresolved and that the primary source recorded even degree four as open. This is supplied dated curation derived from the OWR material, not an independently verified assessment. The record explicitly calls for independent verification of time-sensitive status and cited claims.
\item The registry labels the statement \texttt{corrected\_verified} and reports a manual semantic, typographic, and LaTeX review of its cleanup. These are supplied editorial assertions, not independent mathematical verification of this draft.
\end{itemize}
\end{document}
